\begin{table}[htbp]
\centering
\caption{Representative studies using sensitivity or perturbation-based robustness checks.}
\label{tab:sensitivity_methods}
\small
\begin{tabularx}{\columnwidth}{@{}p{1.8cm}XX@{}}
\toprule
\textbf{Reference} & \textbf{Method} & \textbf{Application} \\
\midrule
\cite{Zhang2021} & Guided saliency + perturbation & Arrhythmia detection \\
\cite{Nurmaini2022} & Grad-CAM sensitivity & Prenatal echo. \\
\cite{Lee2021Cardiomegaly} & Heatmap localisation checks & Cardiomegaly \\
\cite{MorenoSanchez2023} & Attribution stability & HF survival \\
\cite{Chmiel2023} & Saliency sensitivity & Brain tumour \\
\cite{Pertuz2023} & Saliency vs.\ lesion masks & Breast cancer \\
\cite{Repetto2025EvaluatingTR} & Corruption robustness & Medical imaging \\
\cite{Patricio2024} & AOPC + Max-Sensitivity & Dermatol./histopath. \\
\cite{Pandey2024} & qxAI sensitivity metrics & Cardiac disease \\
\cite{arun2021assessing} & Localisation evaluation & Radiology \\
\cite{saporta2022benchmarking} & Saliency benchmarking & CXR interpretation \\
\cite{wollek2023attention} & Attention vs.\ gradient & Pneumothorax \\
\cite{sun2023improving} & Patch perturbation & COVID-19 CXR \\
\cite{cerekci2024quantitative} & Quantitative saliency & Mammography \\
\cite{qin2025consistency} & Consistency benchmarking & Brain MRI \\
\cite{brima2024saliency} & Statistical saliency analysis & Medical imaging \\
\cite{jacob2025indian} & Cross-corpus perturbation & Wound healing \\
\bottomrule
\end{tabularx}
\end{table}
